% Claim proof file for row claim_3_6 -- "SplitTopologicalOrder".
% Auto-generated stub by scaffold/solve_chapter.py at row start. A manager
% agent dedicated to writing the TeX proof fills in the proof body below.
%
% This file is a sibling of `main.tex` inside `<subsection>/tex/`. The
% `subfiles` package lets it render both standalone and as part of
% `main.tex` via `\subfile{...}`.
%
% The statement is restated above the proof so the file is self-contained
% when rendered on its own. The proof must:
%   - Stay close to the lecture notes' proof if one exists (always search
%     the LN first for a `\begin{proof}` block immediately following the
%     claim; copy / lightly edit when present).
%   - Otherwise use the LN paradigm and cite prior claims by their `ref`.
%   - Be self-contained enough that a mathematician can follow it without
%     re-reading the LN.
%   - Have no `sorry`, `omitted`, or "obvious" shortcuts.

\documentclass[main]{subfiles}

\begin{document}
% ref: claim_3_6 (proof, with statement restated for readability)
% title: SplitTopologicalOrder
\def\rowref{claim\_3\_6}
\phantomsection\label{claim_3_6}

% --- statement (restated) ---------------------------------------------------
\begin{Rem}[SplitTopologicalOrder]
    For a CADMG $G=(J,V,E,L)$, also $G_{\spl(W)}$ is acyclic.
    If $<$ is any topological order of $G$ given by enumerating all nodes $v \in J \cup V$ via:
    \[ v_1 < v_2 < \cdots < v_n,\]
    then, for instance, 
    a topological order for $G_{\spl(W)}$ can be achieved by assigning for a node $v_j \in W$ with index $j$ the node 
    $v_j^0$ the index $j-\frac{1}{3}$
    and $v_j^1$ the index $j+\frac{1}{3}$, and then ordering all nodes according to their index value.
\end{Rem}

% --- proof ------------------------------------------------------------------
\begin{proof}
% TODO: write the proof body.
\end{proof}

\end{document}
